\documentclass[11pt, a4paper]{article}

% ── Packages ──────────────────────────────────────────────────────────
\usepackage[utf8]{inputenc}
\usepackage[T1]{fontenc}
\usepackage{lmodern}
\usepackage[margin=1in]{geometry}
\usepackage{amsmath, amssymb, amsthm}
\usepackage{booktabs}
\usepackage{longtable}
\usepackage{graphicx}
\usepackage{hyperref}
\usepackage{xcolor}
\usepackage{listings}
\usepackage{tcolorbox}
\usepackage{polya}

\hypersetup{
  colorlinks=true,
  linkcolor=polyablue,
  urlcolor=polyablue,
  citecolor=polyablue,
}

% ── Code listing style ────────────────────────────────────────────────
\lstdefinestyle{polya-python}{
  language=Python,
  basicstyle=\ttfamily\small,
  keywordstyle=\color{polyablue}\bfseries,
  commentstyle=\color{polyagray}\itshape,
  stringstyle=\color{polyagreen},
  numbers=left,
  numberstyle=\tiny\color{polyagray},
  numbersep=8pt,
  frame=single,
  framerule=0.4pt,
  rulecolor=\color{polyagray},
  breaklines=true,
  showstringspaces=false,
  tabsize=4,
}
\lstset{style=polya-python}
\title{Bin Packing -- Container Loading}
\author{uber-polya}
\date{2026-02-26}

\begin{document}

\maketitle
\tableofcontents
\newpage

% ── Phase A: Problem Formulation ─────────────────────────────────────
\section{Problem Formulation}

\begin{polyamodel}

\textbf{Problem Type}: Find \hfill
\textbf{Domain}: Integer Linear Programming


\end{polyamodel}

% ── Universe ──────────────────────────────────────────────────────────
\subsection{Universe}

\begin{itemize}
  \item $See problem description$
\end{itemize}

% ── Variables ─────────────────────────────────────────────────────────
\subsection{Variables}

\begin{longtable}{lll}
\toprule
\textbf{Symbol} & \textbf{Meaning} & \textbf{Type / Range} \\
\midrule
\endhead
$x$ & decision variable & $\mathbb{{R}}$ \\
\bottomrule
\end{longtable}

% ── Structure ─────────────────────────────────────────────────────────
\subsection{Structure}

- **Items**: 25 packages with weights ranging from 5 to 45 kg (generated from seed=42)
- **Container capacity**: 100 kg each
- **Objective**: Minimize the number of containers used
- **Approaches**: First Fit Decreasing (heuristic) vs. ILP (exact optimal)

% ── Mapping ───────────────────────────────────────────────────────────
\subsection{Real-World Mapping}

\begin{longtable}{ll}
\toprule
\textbf{Real-World Concept} & \textbf{Mathematical Object} \\
\midrule
\endhead
problem input & $instance parameters$ \\
\bottomrule
\end{longtable}

% ── Constraints ───────────────────────────────────────────────────────
\subsection{Constraints}

\begin{enumerate}
  \item $Bin packing**: Classic NP-hard combinatorial optimization problem$
  \hfill \textit{()}
  \item $First Fit Decreasing (FFD)**: Sort items largest-first, place each in the first bin that fits$
  \hfill \textit{()}
  \item $NP-hardness**: No known polynomial-time exact algorithm; heuristics provide fast approximations$
  \hfill \textit{()}
  \item $Approximation ratio**: FFD uses at most ceil(11/9 * OPT + 6/9) bins$
  \hfill \textit{()}
  \item $ILP relaxation**: Integer Linear Programming finds the provably optimal solution$
  \hfill \textit{()}
\end{enumerate}

% ── Objective / Claim ─────────────────────────────────────────────────
\subsection{Claim}


% ── Approach ──────────────────────────────────────────────────────────
\subsection{Approach}

\begin{description}
  \item[Suggested approach:] FFD heuristic + ILP (PuLP/CBC, Branch \& Bound)
  \item[Complexity class:] 
  \item[Available tools:] FFD heuristic + ILP
\end{description}
% ── Phase B: Solution ────────────────────────────────────────────────
\section{Solution}

\begin{polyaresult}

\textbf{Answer}: Solution computed


\textbf{Optimal}: Yes \hfill
\textbf{Feasible}: Yes
\textbf{Algorithm}: FFD heuristic + ILP (PuLP/CBC, Branch \& Bound) \quad $$

\textbf{Time}: 0.029225589998532087s

\textbf{Certificate}: Branch-and-bound solved to optimality.

\end{polyaresult}

% ── Solution Details ──────────────────────────────────────────────────
\subsection{Solution Details}

Solution computed successfully.

% ── Verification ──────────────────────────────────────────────────────
\subsection{Verification}

\begin{longtable}{lcl}
\toprule
\textbf{Check} & \textbf{Status} & \textbf{Detail} \\
\midrule
\endhead
Feasibility & \passmark & All constraints satisfied \\
Optimality & \passmark & FFD heuristic + ILP (PuLP/CBC, Branch \& Bound) \\
\bottomrule
\end{longtable}

% ── Phase C: Interpretation ──────────────────────────────────────────
\section{Interpretation}

% ── The Question & The Answer ─────────────────────────────────────────
\subsection{The Question}
- **Items**: 25 packages with weights ranging from 5 to 45 kg (generated from seed=42)
- **Container capacity**: 100 kg each
- **Objective**: Minimize the number of containers used
- **Approaches**: First Fit Decreasing (heuristic) vs.

\subsection{The Answer}

\begin{polyainsight}[title={Bottom Line}]
A optimal solution was found.
\end{polyainsight}

% ── What This Means ───────────────────────────────────────────────────
\subsection{What This Means}

- FFD heuristic: 7 bins, average utilization \textasciitilde{}93\%
- ILP optimal: 7 bins, average utilization \textasciitilde{}93\%
- Lower bound: ceil(648 / 100) = 7 bins
- Utilization analysis: Both FFD and ILP achieve the lower bound; FFD minimum bin utilization is 59\% while ILP achieves at least 83\%
- All 8 verification checks pass

% ── Sensitivity Analysis ──────────────────────────────────────────────

% ── Figures ───────────────────────────────────────────────────────────

% ── Recommendations ───────────────────────────────────────────────────
\subsection{Recommendations}

\begin{enumerate}
  \item The solution is provably optimal; no further improvement is possible within this model.
\end{enumerate}

% ── Limitations ───────────────────────────────────────────────────────
\subsection{Limitations}

\begin{polyawarnbox}
\begin{itemize}
  \item Model assumes deterministic parameters; real-world variability not captured.
  \item Solution quality depends on the accuracy of input data.
\end{itemize}
\end{polyawarnbox}

% ── Appendix ─────────────────────────────────────────────────────────

\end{document}